% =============================================================
%  ch10_betriebssystem.tex  --  Betriebssystem und Distribution
% =============================================================

\chapter{Betriebssystem und Distribution}

\section{Distributionswahl (Stand Juni 2026)}
ROS\,2-Releases erscheinen jährlich; relevant sind die LTS-Versionen. Für ein
langfristiges Lernprojekt:
\begin{longtable}{@{}p{3.4cm}p{2.6cm}p{2.4cm}p{6cm}@{}}
\toprule
\textbf{Distribution} & \textbf{Ubuntu} & \textbf{Support bis} & \textbf{Eignung} \\
\midrule
\textbf{Jazzy Jalisco} & 24.04 & 2029 (LTS) & \textbf{Empfehlung}: reif, viele Tutorials, von Isaac Sim unterstützt \\
Humble Hawksbill       & 22.04 & 2027 (LTS) & meiste Tutorials, etwas älter \\
Kilted Kaiju           & 24.04 & $\sim$2026 & non-LTS, kurzer Support -- meiden \\
Lyrical Luth           & 24.04 & 2031 (LTS) & brandneu (Mai 2026), Ökosystem noch nicht nachgezogen \\
\bottomrule
\end{longtable}

\begin{infobox}[title=Empfehlung]
\textbf{Ubuntu 24.04 LTS + ROS\,2 Jazzy Jalisco.} LTS bis 2029, ausgereifte
Tutorials, von Gazebo Harmonic und NVIDIA Isaac Sim unterstützt. Lyrical Luth
ist erst $\sim$10 Tage alt -- für den Einstieg zu frisch.
\end{infobox}

\section{VM, Dual-Boot oder WSL2?}
\begin{warnbox}[title=Wichtige Einschränkung bei der VM]
Eine VM (VirtualBox/VMware) ist für die \emph{ROS\,2-Grundlagen} in Ordnung,
aber \textbf{Gazebo Harmonic und Isaac Sim brauchen GPU-Beschleunigung}, die in
VMs schlecht durchgereicht wird -- 3D-Simulation läuft dort zäh. Optionen,
nach Eignung:
\begin{enumerate}
\item \textbf{Natives Ubuntu (Dual-Boot)} -- beste Leistung, empfohlen sobald du
  simulierst.
\item \textbf{WSL2 unter Windows} -- inzwischen mit GPU-Zugriff (WSLg), guter
  Kompromiss.
\item \textbf{VM} -- nur für CLI/Tutorials; für schwere Simulation ungeeignet.
\end{enumerate}
\end{warnbox}
